\subsubsection{A PDE without Solutions}\label{subsec:pde-without-solutions}

Every now and then a paper appears with a title like "A method to solve all partial differential equations." The content of such papers is always very far from satisfying the claims made in the title. It is rumored that a paper of this kind inspired Lewy to construct his famous example of a linear PDE which has no solutions at all. This example also highlights the importance of analyticity in the Cauchy-Kovalevskaya result.

\begin{Theorem}
\label{thm:lewy-nonsolvability}
For a complex-valued function $u(x,y,z)$, let
\begin{equation}\label{eq:lewy-differential-operator}
Lu=-u_x-iu_y+2i(x+iy)u_z.
\end{equation}
Then there is a real-valued function $f(x,y,z)$, of class $C^\infty(\mathbb{R}^3)$, such that the equation
\begin{equation}\label{eq:lewy-inhomogeneous-equation}
Lu=f(x,y,z)
\end{equation}
has no solutions of class $C^1(\Omega)$ in any open subset $\Omega\subset\mathbb{R}^3$.
\end{Theorem}

We note that when $f$ is analytic, the Cauchy-Kovalevskaya theorem applies and noncharacteristic initial-value problems for \eqref{eq:lewy-inhomogeneous-equation} have local solutions. In contrast, for nonanalytic $f$ there may be no solutions, even if no initial conditions are prescribed.

We shall not give a full proof of the theorem, but outline some of the main ideas. First, we shall prove the following lemma.

\begin{Lemma}
\label{lem:lewy-local-nonsolvability}
Let $\psi\in C^\infty(\mathbb{R})$ be real-valued and such that $\psi$ is not real analytic at $z_0$. Then the equation
\begin{equation}\label{eq:lewy-local-obstruction-equation}
Lu=\psi'(z)
\end{equation}
has no solution of class $C^1$ in any neighborhood of the point $(0,0,z_0)$.
\end{Lemma}

\begin{Proof}
Assume the contrary and let $u$ be a solution in a neighborhood of $(0,0,z_0)$, say for $x^2+y^2<\epsilon$, $|z-z_0|<\epsilon$. We set
\begin{equation}\label{eq:lewy-polar-transformed-function}
v(r,\theta,z)=e^{i\theta}\sqrt r\,u(\sqrt r\cos\theta,\sqrt r\sin\theta,z).
\end{equation}
After some algebra, we find that $v$ satisfies the equation
\begin{equation}\label{eq:lewy-polar-differential-equation}
-2v_r-\dfrac irv_\theta+2iv_z=\psi'(z).
\end{equation}
We set
\begin{equation}\label{eq:lewy-angular-average}
V(z,r)=\int_0^{2\pi}v(r,\theta,z)\,\mathrm{d}\theta,
\end{equation}
and integrate \eqref{eq:lewy-polar-differential-equation} with respect to $\theta$. This yields
\begin{equation}\label{eq:lewy-averaged-differential-equation}
V_z+iV_r=-\pi i\psi'(z).
\end{equation}
Clearly $V$ is of class $C^1$ for $0<r<\epsilon$, $|z-z_0|<\epsilon$ and continuous for $0\leq r<\epsilon$, $|z-z_0|<\epsilon$. Moreover, $V(z,0)=0$. Consider now $W=V(z,r)+\pi i\psi(z)$. Then \eqref{eq:lewy-averaged-differential-equation} yields $W_z+iW_r=0$, which makes $W$ a holomorphic function of the complex variable $z+ir$ in the domain $|z-z_0|<\epsilon$, $0<r<\epsilon$. Moreover, $W$ is continuous up to the boundary $r=0$ and it is purely imaginary there. By the Schwarz reflection principle (see Problem~\ref{ex:schwarz-reflection-principle}), we can extend $W$ to a holomorphic function on $|r|<\epsilon$, $|z-z_0|<\epsilon$ by defining $W(z,-r)=-\overline{W(z,r)}$. But this implies that $\pi\psi(z)$, the imaginary part of $W(z,0)$, is real analytic.
\end{Proof}

The equation
\begin{equation}\label{eq:lewy-shifted-forcing}
Lu=\psi'(z-2y_0x+2x_0y)
\end{equation}
is transformed into \eqref{eq:lewy-local-obstruction-equation} by the simple substitution
\begin{equation}\label{eq:lewy-coordinate-translation}
U(x,y,z)=u(x+x_0,y+y_0,z+2y_0x-2x_0y).
\end{equation}
Hence, given any point $(x_0,y_0,z_0)$, we can find a function $f(x,y,z)$ such that \eqref{eq:lewy-inhomogeneous-equation} has no solutions in any neighborhood of $(x_0,y_0,z_0)$.

Consider now a finite number of points $(x_i,y_i,z_i)$, $i=1,\ldots,k$. We shall construct a function $f$ such that \eqref{eq:lewy-inhomogeneous-equation} has no solutions in the neighborhood of any of these points. To do so, we choose $\psi$ to be a real-valued function in $C^\infty(\mathbb{R})$ which is not analytic anywhere (see Problem~\ref{ex:smooth-nowhere-analytic-function}) and make the ansatz
\begin{equation}\label{eq:lewy-finitely-many-forcing-centres}
f(x,y,z)=\sum_{i=1}^kc_i\psi'(z-2y_ix+2x_iy),
\end{equation}
with real coefficients $c_i$. Then, for any choice of $(c_2,c_3,\ldots,c_k)$, there is at most one value $c_1=c_1(c_2,\ldots,c_k)$ for which \eqref{eq:lewy-inhomogeneous-equation} has solutions in any neighborhood of $(x_1,y_1,z_1)$. Assume the contrary, i.e., that there are two such values $c_1$ and $\tilde c_1$. Then we conclude that the equation
\begin{equation}\label{eq:lewy-two-coefficient-contradiction}
Lu=(c_1-\tilde c_1)\psi'(z-2y_1x+2x_1y)
\end{equation}
has a solution in a neighborhood of $(x_1,y_1,z_1)$, in contradiction to the result proved above.

Likewise, for given $(c_1,c_3,\ldots,c_k)$ there is at most one value $c_2=c_2(c_1,c_3,\ldots,c_k)$ for which \eqref{eq:lewy-inhomogeneous-equation} can have a solution in any neighborhood of $(x_2,y_2,z_2)$. Now restrict $c_i$ to a set of $l$ elements, $l>k$. There are $l^k$ possible choices of $(c_1,\ldots,c_k)$. However, there are at most $kl^{k-1}$ choices for which any of the relations $c_1=c_1(c_2,\ldots,c_k)$, $c_2=c_2(c_1,c_3,\ldots,c_k)$, etc. hold. Hence there are choices for which \eqref{eq:lewy-inhomogeneous-equation} has no solutions in any neighborhood of any of the points $(x_i,y_i,z_i)$. In fact, if $l$ is large relative to $k$, this will be the case for "most" choices of the $c_i$.

To complete the proof of the theorem, one needs to extend this argument to a countable number of points. Let $(x_i,y_i,z_i)$, $i\in\mathbb{N}$, be a sequence of points that is dense in $\mathbb{R}^3$. Then we make the ansatz
\begin{equation}\label{eq:lewy-countably-many-forcing-centres}
f(x,y,z)=\sum_{i=1}^{\infty}c_i\psi'(z-2y_ix+2x_iy).
\end{equation}
If $c_i$ converges to zero sufficiently rapidly as $i\to\infty$, then $f$ is of class $C^\infty$. It can still be shown that, for "most" choices of the $c_i$, \eqref{eq:lewy-inhomogeneous-equation} has no solutions in any neighborhood of any point $(x_i,y_i,z_i)$; hence it does not have solutions anywhere. Carrying out this argument is not easy and requires the methods of functional analysis. We shall not pursue this point here and instead refer to the literature, see e.g., \cite{Jo}.

\paragraph*{Problems}

\begin{enumerate}[label=\thesection.\arabic*.,ref=\thesection.\arabic*,start=7,series=book-problems]

\item\label{ex:multi-index-geometric-series-identity} If $|x_i|<1$ for $i=1,\ldots,n$, show that
\begin{equation}\label{eq:multivariable-geometric-product}
\sum_\alpha\bm{x}^{\alpha}=\dfrac1{(1-x_1)(1-x_2)\cdots(1-x_n)}.
\end{equation}
Apply $D^\beta$ to both sides and compare. Then set $\bm{x}=(q,q,\ldots,q)$. This will yield \eqref{eq:multi-index-geometric-derivative-sum}.

\item\label{ex:complete-power-series-facts} Fill in the remaining details in the proof of claims 1--3 about power series (proceed analogously as for power series in one variable).

\item\label{ex:taylor-series-real-analyticity} Let the Taylor series
\begin{equation}\label{eq:exercise-local-taylor-analyticity}
\sum_\alpha c_\alpha\bm{x}^{\alpha}
\end{equation}
converge in $D=\{\bm{x}\in\mathbb{R}^n\mid|\bm{x}|<r\}$. Show that the function represented by the Taylor series is real analytic in $D$.

\item\label{ex:analytic-taylor-convex-domain} Let $f$ be real analytic in a convex domain $D$ which contains the origin. Let
\begin{equation}\label{eq:exercise-convex-domain-taylor-series}
\sum_\alpha c_\alpha\bm{x}^{\alpha}
\end{equation}
be the associated Taylor series. Assume the series converges absolutely for some $\bm{x}\in D$. Show that it converges to $f(\bm{x})$. Hint: Consider $f(\lambda\bm{x})$ for $0<\lambda<1$ and pass to the limit $\lambda\to1$.

\item\label{ex:elliptic-cauchy-instability} Find the solution of the initial-value problem
\begin{equation}\label{eq:elliptic-cauchy-high-frequency-data}
u_{xx}+u_{yy}=0,\quad u(x,0)=0,\quad u_y(x,0)=\dfrac{\sin nx}{n}.
\end{equation}
Discuss what happens as $n\to\infty$. Compare with the analogous initial-value problem for $u_{xx}-u_{yy}=0$.

\item\label{ex:nonlinear-analytic-breakdown} Consider the initial-value problem
\begin{equation}\label{eq:nonlinear-cauchy-characteristic-problem}
u_y=uu_x,\qquad u(x,0)=f(x),
\end{equation}
where $f$ is analytic. Show that unless $f$ is monotone decreasing, there cannot be an analytic (or even continuous) solution for all positive values of $y$. Hint: $u$ is constant along characteristics given by $\mathrm{d}x/\mathrm{d}y=-u$. If $f$ is not monotone decreasing, characteristics must intersect.

\item\label{ex:reduced-initial-problem-equivalence} Show in detail that \eqref{eq:first-order-scalar-nonlinear-system} with initial conditions \eqref{eq:first-order-scalar-nonlinear-initial-data} is equivalent to \eqref{eq:differentiated-scalar-nonlinear-equation} with initial conditions \eqref{eq:differentiated-scalar-initial-data}. Is \eqref{eq:first-order-scalar-nonlinear-system} taken by itself equivalent to \eqref{eq:differentiated-scalar-nonlinear-equation}? If not, how do they differ?

\item\label{ex:laplace-system-standard-form} The system \eqref{eq:first-order-laplace-system} is not of the form \eqref{eq:ck-intro-quasilinear-system}. Discuss how an initial-value problem for \eqref{eq:first-order-laplace-system} can be reduced to the standard form \eqref{eq:ck-intro-quasilinear-system}, \eqref{eq:ck-intro-zero-initial-data}.

\item\label{ex:schwarz-reflection-principle} Let $\Omega$ be a domain in $\mathbb{C}$, symmetric with respect to the real axis. Let $f$ be holomorphic in $\Omega\cap\{\operatorname{Im}z>0\}$ and continuous on $\Omega\cap\{\operatorname{Im}z\geq0\}$. Moreover, assume that $f$ takes real values on $\Omega\cap\{\operatorname{Im}z=0\}$. Show that $f$ can be extended to a function that is holomorphic in all of $\Omega$ by setting $f(z)=\overline{f(\bar z)}$ for $\operatorname{Im}z<0$. Hint: Show that $\int_Cf(z)\,\mathrm{d}z=0$ for any closed rectifiable curve $C$ such that $C$ and its interior lie in $\Omega$. It suffices to show this when $C$ is a triangle.

\item\label{ex:surface-descriptions-equivalent} Show that the three definitions of a $C^k$ (or analytic) surface are indeed equivalent.

\item\label{ex:smooth-nowhere-analytic-function} Show that the function
\begin{equation}\label{eq:smooth-nowhere-analytic-series}
f(x)=\sum_{n=1}^{\infty}\dfrac{\cos(n!x)}{(n!)^n}
\end{equation}
is of class $C^\infty$ on $\mathbb{R}$, but is not real analytic anywhere. Hint: Show first that $f$ is not in $C_{M,r}(0)$ for any $M$ and $r$. Next show that $f(x)-f(x+2\pi q)$ is analytic for any rational number $q$.

\end{enumerate}
